\chapter{Relative Hochschild duality and higher-body couplings}
\label{ch:ym-higher-body-couplings}

\index{Hochschild cohomology!relative|textbf}
\index{Yang--Mills!higher-body couplings}
\index{Kodaira--Spencer class!mixed}

The objects of interest are not single algebras but hierarchies of coherent
multiplications, restrictions, and transgressions. When several boundary
or defect sectors interact, ordinary Hochschild theory must be
relativized: the multi-body problem forces a simplicial package whose
nerve organises all pairwise and higher mixed deformations into a single
coherent datum. The relative Hochschild cochain complex
$CC^\bullet_{\mathrm{rel}}(M_{12};M_1,M_2)$, defined as the
homotopy cofibre of the restriction maps, captures the mixed
couplings invisible from either boundary alone; and the higher-body
\v{C}ech complex extends this to an arbitrary finite collection of
boundary conditions, assembled via a simplicial boundary nerve whose
cohomology controls the simultaneous deformability of the configuration.

Single-boundary Hochschild cohomology $HH^\bullet(M_i)$ and its
Serre-dual deformation algebra are the strict starting data; mixed and
higher-body structures are coherent corrections to the fiction that each
boundary deforms independently. The relative duality theorem
(Theorem~\ref{thm:relative-duality-theorem}) is a projection of Verdier
intertwining in $\mathfrak{g}^{\mathrm{mod}}_\cA$; the long exact
sequence it produces is the Mayer--Vietoris sequence of the modular
deformation complex restricted to the boundary configuration. The
principal $\mathcal{W}$-boundary of twisted Yang--Mills provides the
test case, where the mixed Kodaira--Spencer class and its higher-body
\v{C}ech generalisation are computed explicitly and matched against
complementarity. The simplicial boundary package defined here feeds into
the instanton screening formalism of
Chapter~\ref{ch:ym-instanton-screening}; the Novikov completion of the
higher-body complex controls how nonperturbative sectors mix across
multiple boundaries.

\section{Relative Hochschild duality and the mixed Kodaira--Spencer class}
\label{sec:relative-hochschild-mixed-ks}

Fix a smooth projective curve $X$ and a weakly composable pair of transferred $\mathcal{W}$-dual data
\[
M_1 \xrightarrow{\iota_1} M_{12} \xleftarrow{\iota_2} M_2,
\]
where each $M_i$ is a complete filtered minimal $A_\infty$-algebra quasi-isomorphic to the transferred dual of a principal $\mathcal{W}$-boundary condition.

\begin{definition}[Relative Hochschild cochains]
\label{def:relative-hochschild-cochains}
Define the \emph{relative Hochschild cochain complex}
\[
CC^\bullet_{\mathrm{rel}}(M_{12};M_1,M_2)
:=
\operatorname{hocofib}\!\Big(
CC^\bullet(M_1)\oplus CC^\bullet(M_2)
\xrightarrow{\;\iota_{1*}+\iota_{2*}\;}
CC^\bullet(M_{12})
\Big).
\]
Its cohomology groups are denoted
\[
HH^\bullet_{\mathrm{rel}}(M_{12};M_1,M_2)
:= H^\bullet\big(CC^\bullet_{\mathrm{rel}}(M_{12};M_1,M_2)\big).
\]
The degree-zero group is called the \emph{relative derived center}.
\end{definition}

\begin{definition}[Mixed deformation $L_\infty$-algebra]
\label{def:mixed-deformation-linfty-relative}
Let
\[
\mathfrak g_i:=\tau_{\ge 1}CC^\bullet(M_i)[1],
\qquad
\mathfrak g_{12}:=\tau_{\ge 1}CC^\bullet(M_{12})[1],
\]
with the filtered $L_\infty$ structures from the Gerstenhaber bracket. The restriction maps from $\iota_1,\iota_2$ define an $L_\infty$ morphism
\[
r=(r_1,r_2):\mathfrak g_{12}\to \mathfrak g_1\oplus \mathfrak g_2.
\]
Define the \emph{mixed deformation $L_\infty$-algebra}
\[
\mathfrak g_{\mathrm{mix}}:=\operatorname{hofib}(r).
\]
\end{definition}

\begin{theorem}[Relative duality theorem; \ClaimStatusProvedHere]
\label{thm:relative-duality-theorem}
Assume the degreewise complementarity isomorphisms
\[
H^k(\mathfrak g_i)
\cong
HH^{1-k}(M_i)^\vee\otimes \omega_X,
\qquad
H^k(\mathfrak g_{12})
\cong
HH^{1-k}(M_{12})^\vee\otimes \omega_X
\]
for all $k$ under consideration, and assume the relevant cohomology groups are finite-dimensional. Then for every integer $k$ there is a natural isomorphism
\[
H^k(\mathfrak g_{\mathrm{mix}})
\cong
HH^{1-k}_{\mathrm{rel}}(M_{12};M_1,M_2)^\vee\otimes \omega_X.
\]
In particular,
\[
H^1(\mathfrak g_{\mathrm{mix}})
\cong
HH^{0}_{\mathrm{rel}}(M_{12};M_1,M_2)^\vee\otimes \omega_X,
\]
and
\[
H^2(\mathfrak g_{\mathrm{mix}})
\cong
HH^{-1}_{\mathrm{rel}}(M_{12};M_1,M_2)^\vee\otimes \omega_X.
\]
\end{theorem}

\begin{proof}
The homotopy fiber defining $\mathfrak g_{\mathrm{mix}}$ gives a distinguished triangle
\[
\mathfrak g_{\mathrm{mix}} \to \mathfrak g_{12} \xrightarrow{r} \mathfrak g_1\oplus \mathfrak g_2 \to \mathfrak g_{\mathrm{mix}}[1],
\]
hence a long exact sequence in cohomology. The homotopy cofiber defining $CC^\bullet_{\mathrm{rel}}$ gives a distinguished triangle
\[
CC^\bullet(M_1)\oplus CC^\bullet(M_2)
\to CC^\bullet(M_{12})
\to CC^\bullet_{\mathrm{rel}}(M_{12};M_1,M_2)
\to
\big(CC^\bullet(M_1)\oplus CC^\bullet(M_2)\big)[1],
\]
and therefore a long exact sequence in Hochschild cohomology.

The assumed complementarity isomorphisms identify the cohomology of the first triangle, up to the twist by $\omega_X$, with that of the second under the degree shift $k\mapsto 1-k$. Naturality of the restriction maps identifies the boundary morphisms in the two long exact sequences. Dualizing the second and comparing term-by-term with the first, the five-lemma yields the stated isomorphisms for all~$k$.
\end{proof}

\begin{corollary}[Relative tangent-to-center principle; \ClaimStatusProvedHere]
\label{cor:relative-tangent-center-principle}
The space of genuine mixed first-order couplings satisfies
\[
T^1_{X,\mathrm{mix}}(\cB_{12};\cB_1,\cB_2)
:= H^1(\mathfrak g_{\mathrm{mix}})
\cong
HH^{0}_{\mathrm{rel}}(M_{12};M_1,M_2)^\vee\otimes \omega_X.
\]
\end{corollary}

\begin{corollary}[Relative dual unobstructedness; \ClaimStatusProvedHere]
\label{cor:relative-dual-unobstructedness}
If
\[
HH^{-1}_{\mathrm{rel}}(M_{12};M_1,M_2)=0,
\]
then every genuine mixed first-order coupling extends to second order.
\end{corollary}

\begin{proof}
By Theorem~\ref{thm:relative-duality-theorem}, the primary mixed obstruction space is
\[
H^2(\mathfrak g_{\mathrm{mix}})
\cong
HH^{-1}_{\mathrm{rel}}(M_{12};M_1,M_2)^\vee\otimes \omega_X.
\]
The assumed vanishing forces the obstruction space to be zero.
\end{proof}

\begin{theorem}[Relative central formality; \ClaimStatusProvedHere]
\label{thm:relative-central-formality}
Assume
\[
HH^{1}_{\mathrm{rel}}(M_{12};M_1,M_2)=0
\qquad\text{and}\qquad
HH^{-1}_{\mathrm{rel}}(M_{12};M_1,M_2)=0.
\]
Then the formal moduli problem of genuine mixed deformations is equivalent to the formal affine space on
\[
HH^{0}_{\mathrm{rel}}(M_{12};M_1,M_2)^\vee\otimes \omega_X.
\]
Equivalently, every mixed first-order coupling extends uniquely to all orders up to unique equivalence.
\end{theorem}

\begin{proof}
By Theorem~\ref{thm:relative-duality-theorem},
\[
H^0(\mathfrak g_{\mathrm{mix}})
\cong HH^{1}_{\mathrm{rel}}(M_{12};M_1,M_2)^\vee\otimes \omega_X=0,
\]
\[
H^1(\mathfrak g_{\mathrm{mix}})
\cong HH^{0}_{\mathrm{rel}}(M_{12};M_1,M_2)^\vee\otimes \omega_X,
\]
and
\[
H^2(\mathfrak g_{\mathrm{mix}})
\cong HH^{-1}_{\mathrm{rel}}(M_{12};M_1,M_2)^\vee\otimes \omega_X=0.
\]
Transfer the $L_\infty$ structure on $\mathfrak g_{\mathrm{mix}}$ to a minimal model on its cohomology. Degree counting forces every higher bracket on degree-one classes to land in degree two, which vanishes. The minimal model is abelian on the deformation sector, the Maurer--Cartan equation is empty, and the gauge action is trivial since $H^0=0$.
\end{proof}

\begin{definition}[Mixed Kodaira--Spencer class]
\label{def:mixed-kodaira-spencer-class}
Suppose $M_{12}$ carries a complete descending filtration by mixed weight, and write its $A_\infty$ coderivation as
\[
m = m^{(0)} + m^{(1)} + m^{(2)} + \cdots,
\]
where $m^{(r)}$ raises mixed weight by $r$, and $m^{(0)}$ is the decoupled part inherited from $M_1$ and $M_2$.

The class of $m^{(1)}$ in the degree-one relative Hochschild cohomology,
\[
\kappa_{\mathrm{mix}}(M_{12};M_1,M_2)
:= [m^{(1)}]
\in HH^{1}_{\mathrm{rel}}(M_{12};M_1,M_2),
\]
is called the \emph{mixed Kodaira--Spencer class}.
\end{definition}

\begin{lemma}[Well-definedness and invariance; \ClaimStatusProvedHere]
\label{lem:mixed-ks-well-defined}
The mixed Kodaira--Spencer class is independent of the choice of filtered minimal model, up to the canonical identification of relative Hochschild cohomology under filtered $A_\infty$ quasi-isomorphism.
\end{lemma}

\begin{proof}
A filtered $A_\infty$ quasi-isomorphism between two transferred models intertwines their coderivations up to gauge equivalence in the Hochschild dg Lie algebra. The weight-one component of a gauge-equivalent coderivation changes by a relative Hochschild coboundary. Hence the cohomology class of $m^{(1)}$ is invariant.
\end{proof}

\begin{theorem}[First-correction theorem; \ClaimStatusProvedHere]
\label{thm:first-correction-theorem}
Equip $CC^\bullet_{\mathrm{rel}}(M_{12};M_1,M_2)$ with the filtration induced by mixed weight. Then the associated spectral sequence has
\[
E_1^0 \cong HH^0\bigl(\operatorname{gr} CC^\bullet_{\mathrm{rel}}\bigr),
\]
and the first differential is induced by the adjoint action of the mixed Kodaira--Spencer class:
\[
d_1[\phi] = [\kappa_{\mathrm{mix}},\phi].
\]
In particular, on the quadratic-limit term
\[
E_1^0 \cong \widetilde Z_\infty(M_1)\widehat\otimes \widetilde Z_\infty(M_2),
\]
the class of a naive quadratic mixed coupling survives to the next page if and only if it lies in the kernel of
\[
[\kappa_{\mathrm{mix}},-] :
\widetilde Z_\infty(M_1)\widehat\otimes \widetilde Z_\infty(M_2)
\longrightarrow E_1^1.
\]
\end{theorem}

\begin{proof}
On Hochschild cochains, the differential is the Gerstenhaber commutator with the full coderivation:
\[
d = [m,-] = [m^{(0)},-] + [m^{(1)},-] + [m^{(2)},-] + \cdots.
\]
The zeroth differential on the associated graded is therefore $d_0=[m^{(0)},-]$, and the first derived differential is induced by the next term $[m^{(1)},-]$. By Definition~\ref{def:mixed-kodaira-spencer-class}, the cohomology class of $m^{(1)}$ is precisely $\kappa_{\mathrm{mix}}$, so on the $E_1$-page the first differential is the adjoint action of $\kappa_{\mathrm{mix}}$.
\end{proof}

\begin{corollary}[Persistence criterion for the quadratic law; \ClaimStatusProvedHere]
\label{cor:persistence-criterion-quadratic-law}
The quadratic mixed-coupling law survives to first non-quadratic order if and only if
\[
\kappa_{\mathrm{mix}}(M_{12};M_1,M_2)=0.
\]
More generally, if the mixed coderivation components $m^{(1)},\dots,m^{(r)}$ vanish in relative Hochschild cohomology, then the mixed-weight spectral sequence degenerates through $E_{r+1}$.
\end{corollary}

\begin{proof}
Immediate from Theorem~\ref{thm:first-correction-theorem}. For the second statement, if all differential pieces up to mixed weight $r$ are cohomologically trivial, then $d_1,\dots,d_r$ all vanish.
\end{proof}

\begin{theorem}[Universal interaction brackets on mixed couplings; \ClaimStatusProvedHere]
\label{thm:universal-interaction-brackets}
Let
\[
\mathbb T_{\mathrm{mix}}:=H^1(\mathfrak g_{\mathrm{mix}}),
\qquad
\mathrm{Obs}_{\mathrm{mix}}:=H^2(\mathfrak g_{\mathrm{mix}}).
\]
Homotopy transfer endows $H^\bullet(\mathfrak g_{\mathrm{mix}})$ with a minimal $L_\infty$ structure whose restriction to degree one defines multilinear maps
\[
\lambda_n : \operatorname{Sym}^n(\mathbb T_{\mathrm{mix}}) \longrightarrow \mathrm{Obs}_{\mathrm{mix}},
\qquad n\ge 2.
\]
A formal mixed coupling $\alpha\in \mathbb T_{\mathrm{mix}}\widehat\otimes \mathfrak m$ extends to all orders if and only if
\[
\sum_{n\ge 2}\frac{1}{n!}\lambda_n(\alpha,\dots,\alpha)=0.
\]
The first nonzero $\lambda_n$ is a derived invariant of the filtered $A_\infty$ quasi-isomorphism class of the cospan $(M_1\to M_{12}\leftarrow M_2)$.
\end{theorem}

\begin{proof}
This is the Maurer--Cartan description for the minimal $L_\infty$ model of $\mathfrak g_{\mathrm{mix}}$. Since every input has degree~$1$ and $L_\infty$ brackets have degree $2-n$, the output lies in degree~$2 = \mathrm{Obs}_{\mathrm{mix}}$. Invariance under filtered quasi-isomorphism follows from invariance of the minimal model up to $L_\infty$ isomorphism.
\end{proof}


\section{Higher-body \v{C}ech theory for non-quadratic $\mathcal W$-boundaries}
\label{sec:w-boundary-cech-theory}

\begin{convention}[Derived tangent-to-center input]
\label{conv:derived-tangent-center-input}
Throughout this section, let $X$ be a smooth projective curve and let
\[
\mathfrak N=\{M_S\}_{\emptyset\neq S\subseteq I},\qquad I=\{1,\dots,m\},
\]
be a finite collection of transferred $\mathcal W$-dual data indexed by nonempty subsets of $I$.
For each $S\subseteq I$, write $\cB_S$ for the corresponding boundary condition and assume there is a canonical identification
\[
T^1_X(\cB_S)
\xrightarrow{\sim}
HH^0(M_S)^\vee\otimes \omega_X.
\]
Set
\[
Z_S:=HH^0(M_S),\qquad T_S:=T^1_X(\cB_S).
\]
\end{convention}

\begin{definition}[Boundary simplex]
\label{def:w-boundary-simplex}
A \emph{boundary simplex of order $m$} consists of transferred $\mathcal W$-dual data
\[
\mathfrak N=\{M_S\}_{\emptyset\neq S\subseteq I}
\]
with filtered $A_\infty$-morphisms
\[
\rho_{S,T}\colon M_T\longrightarrow M_S
\qquad (\emptyset\neq T\subseteq S\subseteq I)
\]
satisfying $\rho_{S,S}=\id$ and $\rho_{U,S}\circ \rho_{S,T}=\rho_{U,T}$ whenever
$T\subseteq S\subseteq U$.
\end{definition}

\begin{definition}[Derived-center \v{C}ech complex]
\label{def:w-derived-center-cech}
Fix the natural order on $I$. For a boundary simplex $\mathfrak N$, define the cochain complex
\[
C_Z^p(\mathfrak N):=\bigoplus_{|S|=p+1} Z_S,
\qquad 0\le p\le m-1,
\]
with differential $\delta\colon C_Z^p\to C_Z^{p+1}$ given on the component indexed by
$S=\{s_0<\cdots<s_{p+1}\}$ by
\[
(\delta z)_S
:=
\sum_{j=0}^{p+1}(-1)^j (\rho_{S,S\setminus\{s_j\}})_*(z_{S\setminus\{s_j\}}).
\]
We call $C_Z^\bullet(\mathfrak N)$ the \emph{derived-center \v{C}ech complex} of the simplex.
\end{definition}

\begin{definition}[Tangent restriction complex and pure $m$-body couplings]
\label{def:w-tangent-restriction-complex}
Define the cochain complex
\[
C_T^p(\mathfrak N):=\bigoplus_{|S|=m-p} T_S,
\qquad 0\le p\le m-1,
\]
with differential $\partial\colon C_T^p\to C_T^{p+1}$ given on the component indexed by
$T=\{t_0<\cdots<t_{m-p-2}\}$ by
\[
(\partial\lambda)_T
:=
\sum_{i\notin T}(-1)^{\operatorname{pos}(i,T\cup\{i\})-1}
\operatorname{res}_{T\cup\{i\},T}(\lambda_{T\cup\{i\}}),
\]
where $\operatorname{res}_{S,T}\colon T_S\to T_T$ is the restriction map dual to
$(\rho_{S,T})_*\colon Z_T\to Z_S$ under Convention~\ref{conv:derived-tangent-center-input}.

The space of \emph{pure $m$-body first-order couplings} is
\[
T^{1,[m]}_{X,\mathrm{pure}}(\mathfrak N):=H^0\bigl(C_T^\bullet(\mathfrak N)\bigr).
\]
\end{definition}

\begin{lemma}[Homotopy invariance; \ClaimStatusProvedHere]
\label{lem:w-boundary-simplex-homotopy-invariance}
Let $F\colon \mathfrak N\to \mathfrak N'$ be a morphism of boundary simplices which is an
$A_\infty$-quasi-isomorphism on every face and commutes with all face maps. Then the induced map
\[
C_Z^\bullet(\mathfrak N)\longrightarrow C_Z^\bullet(\mathfrak N')
\]
is an isomorphism of cochain complexes. In particular,
\[
H^p\bigl(C_Z^\bullet(\mathfrak N)\bigr)
\cong
H^p\bigl(C_Z^\bullet(\mathfrak N')\bigr)
\]
for all $p$, and the pure coupling spaces are canonically identified.
\end{lemma}

\begin{proof}
Degree-zero Hochschild cohomology is invariant under $A_\infty$-quasi-isomorphism. Hence
$Z_S\cong Z'_S$ for every face $S$, compatibly with the induced maps on $HH^0$. Taking direct sums over all faces of a fixed cardinality yields an isomorphism in each cochain degree, and the compatibility with face maps makes this an isomorphism of complexes. The final statement follows from Theorem~\ref{thm:w-boundary-cech-duality}.
\end{proof}

\begin{theorem}[Higher-body \v{C}ech duality; \ClaimStatusProvedHere]
\label{thm:w-boundary-cech-duality}
For every boundary simplex $\mathfrak N$ of order $m$, there is a canonical isomorphism of cochain complexes
\[
C_T^p(\mathfrak N)
\xrightarrow{\sim}
\Bigl(C_Z^{m-1-p}(\mathfrak N)\Bigr)^\vee\otimes \omega_X,
\qquad 0\le p\le m-1,
\]
under which the differential $\partial$ is the transpose of $\delta$.
Consequently, for every $0\le q\le m-1$,
\[
H^q\bigl(C_T^\bullet(\mathfrak N)\bigr)
\cong
H^{m-1-q}\bigl(C_Z^\bullet(\mathfrak N)\bigr)^\vee\otimes \omega_X.
\]
In particular,
\[
T^{1,[m]}_{X,\mathrm{pure}}(\mathfrak N)
\cong
H^{m-1}\bigl(C_Z^\bullet(\mathfrak N)\bigr)^\vee\otimes \omega_X.
\]
\end{theorem}

\begin{proof}
By Convention~\ref{conv:derived-tangent-center-input}, each summand $T_S$ is canonically
identified with $Z_S^\vee\otimes \omega_X$. Summing over all subsets $S$ of cardinality $m-p$
gives the objectwise isomorphism
\[
C_T^p(\mathfrak N)
\cong
\Bigl(C_Z^{m-1-p}(\mathfrak N)\Bigr)^\vee\otimes \omega_X.
\]
The sign conventions in Definitions~\ref{def:w-derived-center-cech} and
\ref{def:w-tangent-restriction-complex} were chosen so that the restriction map
$\operatorname{res}_{S,T}$ is exactly the dual transpose of $(\rho_{S,T})_*$. Therefore the alternating
sum defining $\partial$ is the transpose of the alternating sum defining $\delta$.
Taking cohomology yields the stated duality.
\end{proof}

\begin{remark}[Interpretation]
\label{rem:w-boundary-cech-interpretation}
The groups $H^q(C_T^\bullet(\mathfrak N))$ form a tower of first-order descent obstructions.
A degree-$q$ cochain assigns couplings to all faces of codimension $q$; the cocycle condition
requires compatibility on codimension-$(q+1)$ faces. The cohomology class measures
the failure of compatible face data to extend one dimension higher.
Theorem~\ref{thm:w-boundary-cech-duality} identifies this tower with the dual \v{C}ech cohomology of
facewise derived centers.
\end{remark}

\begin{definition}[Decoupled simplex]
\label{def:w-decoupled-simplex}
A boundary simplex $\mathfrak N$ is called \emph{decoupled} if there exist complete vector spaces
$\widetilde Z_1,\dots,\widetilde Z_m$ such that, for every nonempty subset $S\subseteq I$,
\[
Z_S\cong \widehat\bigotimes_{i\in S}(\Bbbk\mathbf{1}\oplus \widetilde Z_i),
\]
and under these identifications the maps $(\rho_{S,T})_*$ are obtained by adjoining the unit in every
factor belonging to $S\setminus T$.
\end{definition}

\begin{definition}[Augmented center complex]
\label{def:w-augmented-center-complex}
For a boundary simplex $\mathfrak N$ of order $m$, define the augmented cochain complex
$\widetilde C_Z^\bullet(\mathfrak N)$ by
\[
\widetilde C_Z^{-1}(\mathfrak N):=\Bbbk,
\qquad
\widetilde C_Z^p(\mathfrak N):=C_Z^p(\mathfrak N)\quad (0\le p\le m-1),
\]
with the augmentation map $\Bbbk\to \bigoplus_i Z_{\{i\}}$ given by the direct sum of unit inclusions.
\end{definition}

\begin{theorem}[Tensor-product model for the augmented center complex; \ClaimStatusProvedHere]
\label{thm:w-augmented-center-tensor-model}
Let $\mathfrak N$ be a decoupled boundary simplex, and for each $i\in I$ set
\[
K_i:=\bigl[\Bbbk\xrightarrow{\eta_i} \Bbbk\mathbf{1}\oplus \widetilde Z_i\bigr],
\]
concentrated in cohomological degrees $0$ and $1$, where $\eta_i(1)=\mathbf{1}$.
Then there is a canonical isomorphism of cochain complexes
\[
\widetilde C_Z^\bullet(\mathfrak N)
\xrightarrow{\sim}
\left(\widehat\bigotimes_{i=1}^m K_i\right)[-1].
\]
Consequently,
\[
H^p\bigl(C_Z^\bullet(\mathfrak N)\bigr)=0
\quad (0\le p<m-1),
\qquad
H^{m-1}\bigl(C_Z^\bullet(\mathfrak N)\bigr)
\cong
\widehat\bigotimes_{i=1}^m \widetilde Z_i.
\]
\end{theorem}

\begin{proof}
Expand the completed tensor product $\widehat\bigotimes_i K_i$. A pure tensor lies in cohomological
degree $r$ exactly when $r$ of the factors lie in $\Bbbk\mathbf{1}\oplus \widetilde Z_i$ and the remaining
$m-r$ lie in $\Bbbk$. Such a choice is equivalent to a subset $S\subseteq I$ of cardinality $r$, and the
corresponding summand is canonically
\[
\widehat\bigotimes_{i\in S}(\Bbbk\mathbf{1}\oplus \widetilde Z_i)
\cong Z_S.
\]
Thus the degree-$r$ part of $\widehat\bigotimes_i K_i$ identifies with
\[
\bigoplus_{|S|=r} Z_S = C_Z^{r-1}(\mathfrak N),
\]
which explains the shift by $[-1]$.

Under this identification, the tensor-product differential inserts the unit in one additional factor, with
the usual Koszul sign determined by the position of that factor. This is exactly the alternating sum in the
augmented \v{C}ech differential of Definition~\ref{def:w-augmented-center-complex}. Hence the two complexes
are canonically isomorphic.

Each $K_i$ has $H^0(K_i)=0$ and $H^1(K_i)\cong \widetilde Z_i$. The completed K\"unneth theorem therefore gives
\[
H^r\left(\widehat\bigotimes_{i=1}^m K_i\right)=0\quad (r\neq m),
\qquad
H^m\left(\widehat\bigotimes_{i=1}^m K_i\right)
\cong
\widehat\bigotimes_{i=1}^m \widetilde Z_i.
\]
After shifting by $[-1]$, the cohomology of $\widetilde C_Z^\bullet(\mathfrak N)$ is concentrated in degree $m-1$.
Since $\widetilde C_Z^\bullet$ and $C_Z^\bullet$ agree in nonnegative degrees, the same statement holds for
$C_Z^\bullet$.
\end{proof}

\begin{corollary}[Pure $m$-body coupling theorem; \ClaimStatusProvedHere]
\label{cor:w-pure-mbody-coupling-theorem}
Let $\mathfrak N$ be a decoupled boundary simplex of order $m$. Then
\[
T^{1,[m]}_{X,\mathrm{pure}}(\mathfrak N)
\cong
\left(\widehat\bigotimes_{i=1}^m \widetilde Z_i\right)^\vee\otimes \omega_X.
\]
Pure $m$-body first-order couplings are thus determined by the full tensor product of the reduced
facewise derived centers.
\end{corollary}

\begin{proof}
Combine Theorem~\ref{thm:w-boundary-cech-duality} with
Theorem~\ref{thm:w-augmented-center-tensor-model}.
\end{proof}

\begin{theorem}[Higher Mayer--Vietoris descent for couplings; \ClaimStatusProvedHere]
\label{thm:w-higher-mayer-vietoris-couplings}
Let $\mathfrak N$ be any boundary simplex of order $m$.
\begin{enumerate}[label=(\roman*)]
\item The obstruction to descending a compatible family of codimension-$1$ couplings
\[
\{\lambda_S\in T_S\}_{|S|=m-1}
\]
to a coupling on the top face $I$ is the cohomology class
\[
[\lambda]\in H^1\bigl(C_T^\bullet(\mathfrak N)\bigr)
\cong
H^{m-2}\bigl(C_Z^\bullet(\mathfrak N)\bigr)^\vee\otimes \omega_X.
\]
\item If $H^{m-2}(C_Z^\bullet(\mathfrak N))=0$, then every compatible codimension-$1$ family descends to
some coupling on the top face.
\item If, in addition, $H^{m-1}(C_Z^\bullet(\mathfrak N))=0$, then such a descent is unique.
\end{enumerate}
More generally, $H^q(C_T^\bullet(\mathfrak N))$ is the obstruction space for descending compatible
codimension-$q$ face data one step upward.
\end{theorem}

\begin{proof}
A codimension-$1$ family is by definition a degree-$1$ cochain in $C_T^\bullet(\mathfrak N)$. It is
compatible on codimension-$2$ faces exactly when it is a cocycle. Such a cocycle descends from a top-face
coupling precisely when it lies in the image of
\[
\partial\colon C_T^0(\mathfrak N)\to C_T^1(\mathfrak N),
\]
so the obstruction is its class in $H^1(C_T^\bullet)$. This proves (i), and (ii) is the vanishing criterion
$H^1(C_T^\bullet)=0$ rewritten via Theorem~\ref{thm:w-boundary-cech-duality}. If a descent exists, any two
choices differ by a degree-$0$ cocycle, i.e.\ by an element of $H^0(C_T^\bullet)$; thus the descent is unique
when $H^0(C_T^\bullet)=0$, which is equivalent to $H^{m-1}(C_Z^\bullet)=0$ by Theorem~\ref{thm:w-boundary-cech-duality}.
The final statement is the same argument in degree~$q$.
\end{proof}

\begin{definition}[Filtered boundary simplex]
\label{def:w-filtered-boundary-simplex}
A boundary simplex $\mathfrak N$ is \emph{filtered} if each $Z_S$ carries a complete descending filtration
$F^\bullet Z_S$ such that every face map $(\rho_{S,T})_*$ preserves filtration.
\end{definition}

\begin{theorem}[Filtered stability of pure $m$-body couplings; \ClaimStatusProvedHere]
\label{thm:w-filtered-stability-pure-mbody}
Let $\mathfrak N$ be a filtered boundary simplex. Assume its associated graded simplex
$\operatorname{gr}\mathfrak N$ is decoupled in the sense of Definition~\ref{def:w-decoupled-simplex}, with reduced
centers $\widetilde Z_1,\dots,\widetilde Z_m$. Then:
\begin{enumerate}[label=(\alph*)]
\item the filtered complex $C_Z^\bullet(\mathfrak N)$ has a convergent spectral sequence whose $E_1$-page is
\[
E_1^p \cong H^p\bigl(C_Z^\bullet(\operatorname{gr}\mathfrak N)\bigr);
\]
\item this spectral sequence collapses at $E_1$;
\item the top cohomology inherits a canonical filtration with
\[
\operatorname{gr}H^{m-1}\bigl(C_Z^\bullet(\mathfrak N)\bigr)
\cong
\widehat\bigotimes_{i=1}^m \widetilde Z_i,
\]
while $H^p(C_Z^\bullet(\mathfrak N))=0$ for $0\le p<m-1$;
\item therefore the pure $m$-body coupling space carries a canonical filtration satisfying
\[
\operatorname{gr}T^{1,[m]}_{X,\mathrm{pure}}(\mathfrak N)
\cong
\left(\widehat\bigotimes_{i=1}^m \widetilde Z_i\right)^\vee\otimes \omega_X.
\]
\end{enumerate}
\end{theorem}

\begin{proof}
The facewise filtrations on the $Z_S$ induce a complete filtration on the direct sums
$C_Z^p(\mathfrak N)$, hence on the whole \v{C}ech complex. Filtered-complex theory gives a convergent
spectral sequence with
\[
E_1^p\cong H^p\bigl(C_Z^\bullet(\operatorname{gr}\mathfrak N)\bigr).
\]
By assumption, the associated graded simplex is decoupled, so Theorem~\ref{thm:w-augmented-center-tensor-model}
shows that $E_1^p=0$ for $p<m-1$ and
\[
E_1^{m-1}\cong \widehat\bigotimes_{i=1}^m \widetilde Z_i.
\]
Since the $E_1$-page is concentrated in a single cohomological degree, all higher differentials vanish and the
spectral sequence collapses at $E_1$. Therefore the final cohomology is also concentrated in degree $m-1$, and its
associated graded is the stated tensor product. The final assertion follows from Theorem~\ref{thm:w-boundary-cech-duality}.
\end{proof}

\begin{remark}[Filtered stability]
\label{rem:w-filtered-stability-meaning}
Once the facewise derived centers admit a compatible filtration with decoupled associated
graded, the first-order pure $m$-body coupling theory is a \emph{filtered deformation} of the quadratic tensor-product law.
The non-quadratic $\mathcal W$-sector does not destroy pure higher-body couplings at first order; it refines them.
\end{remark}
